\subsection*{Notas de la sección}

La explicación del paradigma está basada en \textcite{cormen}. 

En merge sort, la implementación de los pasos de dividir y conquistar es la de \textcite{cormen}, \textcite{skiena} y \textcite{cardoso}. El paso de combinar no es idéntico al de alguna fuente, es similar al de \textcite{cormen2022}, aunque allí usan una variable de más que le suma cuatro instrucciones al algoritmo y, a mi juicio, le resta claridad.

La sección de máximo subarreglo toma como fuente principal \textcite{bentley}, aunque usa la implementación de dividir y conquistar de \textcite{cormen} porque muestra el paradigma más explícitamente. El artículo original de Bentley es muy bueno e incluye también la solución por programación dinámica (que se expone en la siguiente sección).
